\section{Related Works}

\subsection{Cache Eviction Methods}

In the long-sequence inference, the vast scale of the KV cache elements leads to a memory-bound situation, causing significant memory burden and I/O latency~\cite{wang2023catalyst}.
Numerous studies have sought to mitigate by reducing the cache size, notably through the eviction of non-critical cache elements\footnote{For additional related works, see Appendix \ref{apdx:additional_related}}.
These  methods are primarily divided into two categories: \textit{sliding window eviction} and \textit{Top-k eviction} methods.
The sliding window eviction methods~\cite{beltagy2020longformer,han2024lminfinitezeroshotextremelength,SLM}, exemplified by \textit{StreamingLLM}~\cite{SLM}, simply retain several initial cache elements and those within a sliding window, while evicting others. However, the undiscriminating sliding eviction of cache elements results in a significant reduction in generation quality.
In contrast, Top-k eviction methods~\cite{FastGen,liu2024scissorhands,ren2024efficacy,H2o,PyramidInfer,PyramidKV,SnapKV} identify and retain a selected set of $k$ critical cache elements based on attention weights, for enhancing the post-eviction generation quality.
The adaptive budget allocation presented in this paper, tailored for Top-$k$ Eviction Methods, enhances post-eviction generation quality within the same overall budget.
	

In the study of Top-$k$ eviction methods, early work like FastGen~\cite{FastGen} searches and combines multiple strategies, such as maintaining caches of special elements, punctuation elements, recent elements, and Top-$k$ selected elements, based on the characteristics of attention heads.
H2O, as the representation~\cite{H2o,liu2024scissorhands,ren2024efficacy}, develops a Top-$k$ based eviction scheme that leverages the query states of all tokens to identify critical cache elements.
However, under unidirectional mask in LLM computations, aggregating attention weights across all query states often causes recent KV cache elements to be mistakenly evicted, degrading the quality of subsequent generations. 
Recent works, such as SnapKV~\cite{SnapKV} address this issue by using query states within an observation window to identify critical elements, thereby achieving SOTA performance. Subsequent methods, such as Pyramid~\cite{PyramidInfer,PyramidKV}, further introduce budget allocation across different layers.
However, existing Top-$k$ eviction methods typically assign the overall budget uniformly across different heads.
In contrast, Ada-KV, which has demonstrated superior theoretical and empirical results through adaptive budget allocation, provides a novel strategy to optimizing these methods.

\subsection{Sparse Attention Methods}
Sparse attention methods are conceptually related to KV cache eviction, but differ fundamentally in their approach~\cite{jiang2024minference,li2025mminference,liu2024retrievalattention,chen2024arkvale,zhang2025pqcache,sun2025breaking}. The key difference is that KV cache eviction retains only a subset of the KV cache, while sparse attention methods keep all entries but selectively utilize only a critical subset during computation. Consequently, sparse attention methods do not reduce the memory footprint of the KV cache, thus often necessitating cache offloading to CPU memory~\cite{liu2024retrievalattention,chen2024arkvale,zhang2025pqcache} or even disk storage~\cite{sun2025breaking}. Additionally, the two technique lines are, in fact, orthogonal. Future research could explore first employing KV cache eviction to compress the cache to a certain proportion and then applying sparse attention for further acceleration. Ada-KV presented in this paper leverages head-wise adaptive allocation to enhance cache eviction methods, and this approach can similarly be applied in future work to enhance dynamic sparse attention methods.
